\chapter{已有科研基础与所需的科研条件}

\section{基于几何先验引导的表征学习}

目前，该部分研究内容已完成，相关成果已发表、被接收 \citep{peng2025synthetic,peng2026visl}；下面展示相应的消融实验以及和主流方法的对比结果来说明所提出方法的有效性，采用的评价指标为Top-1准确率。

如表\ref{tab:condition_geo_sota}所示，ViSL在仅使用骨架模态的情况下取得了67.14\%的测试准确率，显著高于已有的基于RGB、关键点和多模态融合方法。与最强的融合方法相比，本研究方法提升9.17\%准确率；与骨架模态中表现最好的方法相比，提升21.15\%准确率。这一结果证明了引入几何先验、同时利用不同姿态估计器互补信息的有效性。同时，由于推理阶段仅保留主分支，方法不需要引入多模态融合或模型集成，具备进一步部署到真实系统中的轻量化潜力。

\begin{table}[!htbp]
    \caption{基于几何先验方法与主流方法在MM-WLAuslan上的性能对比（Top-1，\%）。其中，$\dagger$表示由\cite{shen2024mm}实现的结果，$\ddagger$表示由\cite{shen2025cvis}实现的结果，$*$表示额外融合深度模态。}
    \label{tab:condition_geo_sota}
    \setlength{\tabcolsep}{8.5pt}
    \centering
    \begin{tabular}{@{}lcccl@{}}
        \toprule
        \multicolumn{1}{c|}{\multirow{2}{*}{方法}} & \multicolumn{3}{c|}{模态} & \multicolumn{1}{c}{\multirow{2}{*}{Test}} \\ \cmidrule(lr){2-4}
        \multicolumn{1}{c|}{} & RGB & 关键点 & \multicolumn{1}{l|}{融合} & \multicolumn{1}{c}{} \\ \midrule
        \multicolumn{5}{c}{CV-ISLR Challenge~\cite{shen2025cross}方法} \\ \midrule
        \multicolumn{1}{l|}{gkdx2~\cite{wang2025exploiting}} & \checkmark & & \multicolumn{1}{l|}{} & 20.29 \\
        \multicolumn{1}{l|}{WANGXINYU1~\cite{wang2025cross}} & \checkmark & & \multicolumn{1}{l|}{} & 25.39 \\
        \multicolumn{1}{l|}{VIPL\_SLP\cite{peng2025synthetic}} & \checkmark & & \multicolumn{1}{l|}{} & 36.15 \\
        \multicolumn{1}{l|}{VIPL\_SLP\cite{peng2025synthetic}} & & \checkmark & \multicolumn{1}{l|}{} & 45.99 \\
        \multicolumn{1}{l|}{gkdx2*~\cite{wang2025exploiting}} & & & \multicolumn{1}{l|}{\checkmark} & 24.53 \\
        \multicolumn{1}{l|}{WANGXINYU1*~\cite{wang2025cross}} & & & \multicolumn{1}{l|}{\checkmark} & 33.97 \\
        \multicolumn{1}{l|}{Tonicemerald~\cite{zheng2025zero}} & & & \multicolumn{1}{l|}{\checkmark} & 40.30 \\
        \multicolumn{1}{l|}{VIPL\_SLP\cite{peng2025synthetic}} & & & \multicolumn{1}{l|}{\checkmark} & 56.87 \\
        \multicolumn{1}{l|}{VIPL\_SLP*\cite{peng2025synthetic}} & & & \multicolumn{1}{l|}{\checkmark} & 57.97 \\ \midrule
        \multicolumn{5}{c}{现有SOTA方法} \\ \midrule
        \multicolumn{1}{l|}{$\text{VKNet-V}\dagger$~\cite{zuo2023natural}} & \checkmark & & \multicolumn{1}{l|}{} & \multicolumn{1}{r}{14.53} \\
        \multicolumn{1}{l|}{$\text{STC-SLR}\ddagger$~\cite{zhao2024self}} & & \checkmark & \multicolumn{1}{l|}{} & \multicolumn{1}{r}{25.72} \\
        \multicolumn{1}{l|}{$\text{VKNet-K}\ddagger$~\cite{zuo2023natural}} & & \checkmark & \multicolumn{1}{l|}{} & 28.04 \\
        \multicolumn{1}{l|}{$\text{DSTA-SLR}\ddagger$~\cite{hu2024dynamic}} & & \checkmark & \multicolumn{1}{l|}{} & 28.66 \\
        \multicolumn{1}{l|}{CMVSR~\cite{shen2025cvis}} & & \checkmark & \multicolumn{1}{l|}{} & 43.12 \\
        \multicolumn{1}{l|}{$\text{UMDR*}\dagger$~\cite{zhou2023unified}} & & & \multicolumn{1}{l|}{\checkmark} & 22.82 \\
        \multicolumn{1}{l|}{$\text{NLA-SLR}\ddagger$~\cite{zuo2023natural}} & & & \multicolumn{1}{l|}{\checkmark} & 33.23 \\ \midrule
        \multicolumn{1}{l|}{Baseline (P-Stream)} & & \checkmark & \multicolumn{1}{l|}{} & 45.53 \\
        \multicolumn{1}{l|}{Baseline (L-Stream)} & & \checkmark & \multicolumn{1}{l|}{} & 52.75 \\
        \multicolumn{1}{l|}{Ours} & & \checkmark & \multicolumn{1}{l|}{} & \textbf{67.14} \\ \bottomrule
    \end{tabular}
\end{table}

\FloatBarrier

表\ref{tab:condition_geo_keypoint}进一步验证了不同关键点来源的作用。直接使用二维估计关键点训练和测试时，模型在侧视角验证集和测试集上分别达到50.26\%和45.53\%，说明高质量二维定位对手语动作细节建模十分重要。相比之下，仅依赖三维人体重建后投影得到的合成关键点时，侧视角性能明显下降，反映出三维估计和投影过程会引入关键点误差与格式差异。然而，当在合成关键点上加入视角增广后，侧视角验证性能由24.11\%提升至57.78\%，测试性能由22.60\%提升至52.75\%，表明三维结构所提供的视角旋转能力能够有效扩展训练分布。上述结果共同说明，二维关键点具备更准确的空间定位能力，三维估计则提供了有价值的深度与视角变化信息，二者具有明显互补性。

\begin{table}[H]
    \centering
    \caption{估计关键点与合成关键点的消融实验（Top-1，\%）}
    \label{tab:condition_geo_keypoint}
    \begin{tabular}{lcccc}
        \hline
        训练数据 & 测试数据 & Val(F) & Val(S) & Test \\
        \hline
        Estimated & Estimated & 92.58 & 50.26 & 45.53 \\
        Synthetic & Synthetic & 86.17 & 24.11 & 22.60 \\
        Synthetic (aug) & Synthetic & 86.49 & 57.78 & 52.75 \\
        Synthetic (aug) & Estimated & 70.67 & 49.81 & 45.45 \\
        \hline
    \end{tabular}
\end{table}

表\ref{tab:condition_geo_loss}展示了跨视角对比学习各组成部分的贡献。仅使用主分支P-Stream时，模型在正面验证集上性能较高，但侧视角泛化能力有限；仅使用辅助提升分支L-Stream并引入视角增广后，侧视角性能有所提升，但正面视角性能受到三维估计噪声影响。加入跨视角对比损失$\mathcal{L}_{CV}$后，测试准确率由52.75\%提升至54.94\%，说明直接约束不同视角下同类样本的特征一致性能够增强模型的视角不变性。进一步联合P-Stream与L-Stream后，测试准确率提升至65.31\%，表明精确二维定位与三维深度线索可以互补。最后，引入跨流对比损失$\mathcal{L}_{CS}$和代理辅助对比损失$\mathcal{L}_{PCS}$后，测试性能达到67.14\%，同时正面验证集性能也恢复到92.64\%。这些结果说明，所设计的双流对比学习框架能够在利用三维几何先验的同时缓解估计噪声和关键点格式差异。

\begin{table}[H]
    \centering
    \caption{跨视角对比学习设计的消融实验（Top-1，\%）}
    \label{tab:condition_geo_loss}
    \setlength{\tabcolsep}{4pt}
    \begin{tabular}{ccccc|ccc}
        \hline
        P-Stream & L-Stream & $\mathcal{L}_{CV}$ & $\mathcal{L}_{CS}$ & $\mathcal{L}_{PCS}$ & Val(F) & Val(S) & Test \\
        \hline
        \checkmark &  &  &  &  & 92.58 & 50.26 & 45.53 \\
         & \checkmark &  &  &  & 86.49 & 57.78 & 52.75 \\
         & \checkmark & \checkmark &  &  & 86.94 & 59.69 & 54.94 \\
        \checkmark & \checkmark & \checkmark &  &  & 91.63 & 70.62 & 65.31 \\
        \checkmark & \checkmark & \checkmark & \checkmark &  & 92.02 & 71.44 & 66.13 \\
        \checkmark & \checkmark & \checkmark & \checkmark & \checkmark & \textbf{92.64} & \textbf{71.71} & \textbf{67.14} \\
        \hline
    \end{tabular}
\end{table}

\FloatBarrier

\section{基于语义先验引导的表征学习}

目前，该部分内容已初步验证利用语义先验引导的有效性，以下从单模态性能、跨数据集泛化能力以及传播注意力机制三个方面总结已有科研基础。

表\ref{tab:condition_semantic_single_modal}展示了EASL在MM-WLAuslan跨视角设置下的单模态结果。仅使用RGB模态时，EASL将测试集平均准确率由41.63\%提升至54.55\%，说明姿态引导的生成式增广能够有效改善RGB模型对视角、背景和外观变化的适应能力。

\begin{table}[H]
    \centering
    \caption{单模态EASL的消融实验（Top-1，\%）}
    \label{tab:condition_semantic_single_modal}
    \setlength{\tabcolsep}{3.5pt}
    \begin{tabular}{@{}lccc@{}}
        \toprule
        方法 & Val(F) & Val(S) & Test(S) \\
        \midrule
        RGB Baseline & 92.52 & 44.72 & 41.63 \\
        EASL & 92.05 & 56.91 & \textbf{54.55} \\
        \bottomrule
    \end{tabular}
\end{table}

为进一步验证方法的跨数据集泛化能力，已有实验在NationalCSL-DP\citep{jin2025nationalcsl}上进行了正面视角到左视角的迁移评估。表\ref{tab:condition_semantic_nationalcsl}表明，EASL在Val(L)上由RGB基线的3.66\%提升至12.96\%，在Test(L)上由2.72\%提升至10.07\%；同时，EASL在正面验证集Val(F)上也达到96.41\%，优于RGB基线的94.74\%。这些结果说明，语义先验引导的生成式增广并非只对单一数据集有效，也能够在不同语言和不同采集条件的数据集上提升跨视角迁移能力。

\begin{table}[H]
    \centering
    \caption{NationalCSL-DP上的跨数据集验证结果（Top-1，\%）}
    \label{tab:condition_semantic_nationalcsl}
    \setlength{\tabcolsep}{6pt}
    \begin{tabular}{@{}lccc@{}}
        \toprule
        方法 & Val(F) & Val(L) & Test(L) \\
        \midrule
        RGB Baseline & 94.74 & 3.66 & 2.72 \\
        EASL & 96.41 & 12.96 & 10.07 \\
        \bottomrule
    \end{tabular}
\end{table}

表\ref{tab:condition_semantic_keyframe_strategy}比较了在相同生成帧数预算下，全序列与关键帧合成策略的性能差异。可以看到，在不同生成规模下，关键帧合成均优于直接生成完整序列的方式。例如，当生成帧数为100,000时，关键帧策略将Test(S)准确率由47.77\%提升至50.05\%；当生成帧数降低至8,000时，关键帧策略仍由43.30\%提升至44.48\%。这说明，在相同生成成本下，优先合成具有代表性的关键帧，并将其中的外观和场景线索传播到完整序列，比直接逐帧生成更有利于转化为识别性能增益。同时，随着生成帧数从8,000增加到100,000，关键帧策略的Test(S)准确率由44.48\%逐步提升至50.05\%，表明EASL能持续受益于更大规模的合成增广数据。

\begin{table}[H]
    \centering
    \caption{全序列合成与关键帧合成策略的性能对比（Top-1，\%）。其中“生成帧数”表示生成图像帧的总数。}
    \label{tab:condition_semantic_keyframe_strategy}
    \setlength{\tabcolsep}{3.5pt}
    \begin{tabular}{llccc}
        \toprule
        生成帧数 & 合成策略 & Val(F) & Val(S) & Test(S) \\
        \midrule
        \multirow{2}{*}{100,000} & full length & 92.47 & 50.75 & 47.77 \\
                                 & keyframe & 93.76 & 53.61 & 50.05 \\
        \midrule
        \multirow{2}{*}{80,000} & full length & 92.19 & 48.87 & 46.51 \\
                                & keyframe & 93.51 & 52.53 & 49.00 \\
        \midrule
        \multirow{2}{*}{60,000} & full length & 91.21 & 45.98 & 43.07 \\
                                & keyframe & 93.69 & 51.56 & 48.21 \\
        \midrule
        \multirow{2}{*}{20,000} & full length & 92.77 & 47.36 & 44.04 \\
                                & keyframe & 92.81 & 49.18 & 46.20 \\
        \midrule
        \multirow{2}{*}{8,000} & full length & 91.93 & 45.99 & 43.30 \\
                               & keyframe & 92.57 & 47.79 & 44.48 \\
        \bottomrule
    \end{tabular}
\end{table}

表\ref{tab:condition_semantic_propagation}进一步验证了传播注意力机制的作用。去除传播注意力后，模型在Test(S)上的准确率为41.14\%；加入传播注意力后，性能提升至50.05\%。该结果说明，仅依赖姿态图像或稀疏生成关键帧难以充分建模RGB纹理和场景信息，而传播注意力能够将关键帧中的外观、背景和风格线索传递到完整序列，从而缓解关键帧生成与全序列识别之间的信息缺口。

\begin{table}[H]
    \centering
    \caption{传播注意力机制的消融实验（Top-1，\%）}
    \label{tab:condition_semantic_propagation}
    \setlength{\tabcolsep}{5pt}
    \begin{tabular}{@{}lccc@{}}
        \toprule
        Propagation Attention & Val(F) & Val(S) & Test(S) \\
        \midrule
        w/ & 93.76 & 53.61 & 50.05 \\
        w/o & 91.06 & 44.51 & 41.14 \\
        \bottomrule
    \end{tabular}
\end{table}

\FloatBarrier

\section{基于自主移动机器人的自适应手语理解系统}

基于自主移动机器人的自适应手语理解系统是前述算法研究面向真实应用场景的进一步延伸。前两节已经从几何先验和语义先验两个层面验证了模型在准正面视角、侧视角和复杂背景扰动下的鲁棒性：一方面，基于三维几何结构的视角增广与双流对比学习能够缓解固定正面训练数据带来的跨视角泛化不足问题；另一方面，基于语义先验的生成式增广能够增强RGB表征对背景、光照、外观和场景变化的适应能力。这些已有结果为移动机器人系统提供了较为稳定的视觉编码基础，使系统在机器人主动移动过程中面对视角变化、尺度变化和局部遮挡时，仍有可能保持可靠的手语理解性能。

在系统研究基础方面，本文已围绕可穿戴式系统、固定视觉终端以及公共服务场景中的手语翻译系统开展了文献调研。现有研究表明，可穿戴式方案通常能够获得较稳定的运动信号，但会增加用户佩戴负担，影响手语表达的自然性；固定摄像头方案使用负担较低，却容易受到用户站位、身体朝向和遮挡关系影响。本文拟采用移动机器人作为手语理解系统的主动感知载体，核心目的在于将“用户适应设备”的交互方式转变为“设备主动适应用户”的交互方式。机器人可以在会议、课堂和公共服务咨询等多人交流场景中，根据手语者位置、人体朝向、手部可见性和场景障碍物分布主动调整录制位置与相机朝向，从而为手语翻译模型提供更稳定的输入。

目前，该部分已经形成较完整的系统功能设计。系统整体将由移动机器人平台、视觉感知模块、主动决策模块、手语翻译模块和交互反馈模块组成。其中，视觉感知模块负责完成人体检测、手语者跟踪、人体关键点估计和观测质量评估；主动决策模块根据观测质量、翻译置信度和局部可通行约束规划机器人移动路径，并在不明显打断交流的前提下寻找更适合手语录制的视角；手语翻译模块以前述鲁棒表征学习方法作为编码基础，将采集到的手语视频片段转换为自然语言文本；交互反馈模块则向用户呈现翻译文本、置信度和必要提示，并支持确认、纠错、重新采集等交互过程。通过上述设计，系统不仅关注模型预测准确率，也关注真实交流中的可用性、实时性和用户体验。

在实验条件方面，本研究已经具备开展系统验证的算法与数据基础。几何先验部分在MM-WLAuslan数据集上完成了跨视角手语识别验证，语义先验部分进一步在MM-WLAuslan和NationalCSL-DP数据集上验证了跨视角和跨数据集泛化能力，这些实验可作为机器人系统中视觉编码器和翻译模型选择的离线依据。后续系统实验将进一步搭建室内会议式交互场景，采集不同手语者、不同发言位置、不同拍摄距离、不同光照和不同遮挡条件下的视频，并记录机器人位姿、规划路径、相机视频、关键点序列、翻译文本和交互反馈日志。评价指标将同时覆盖翻译质量、系统实时性、主动观测效果和用户交互体验，以验证移动机器人主动感知、路径规划与实时反馈机制是否能够在真实场景下提升手语理解系统的稳定性和自然交互能力。

总体来看，该部分目前已完成问题定义、文献调研、系统架构设计和实验方案设计，且前述鲁棒手语识别算法已经为系统实现提供了核心模型基础。后续工作将重点完善机器人端视觉感知与主动规划模块，实现手语翻译模型在移动平台上的部署，并通过真实交互实验评估系统在会议等动态场景中的实用效果。
